\documentclass[12pt,a4paper]{article}

\usepackage[utf8]{inputenc}
\usepackage[T1]{fontenc}
\usepackage[french]{babel}
\usepackage{geometry}
\usepackage{graphicx}
\usepackage{hyperref}
\usepackage{fancyhdr}
\usepackage{xcolor}
\usepackage{listings}
\usepackage{array}
\usepackage{booktabs}
\usepackage{float}
\usepackage{caption}

\geometry{top=2.5cm, bottom=2.5cm, left=2.5cm, right=2.5cm}

\hypersetup{
    colorlinks=true,
    linkcolor=blue!70!black,
    urlcolor=blue!70!black,
    citecolor=blue!70!black
}

\definecolor{codebg}{rgb}{0.95,0.95,0.95}
\definecolor{codeframe}{rgb}{0.7,0.7,0.7}

\lstset{
    basicstyle=\ttfamily\small,
    backgroundcolor=\color{codebg},
    frame=single,
    rulecolor=\color{codeframe},
    breaklines=true,
    showstringspaces=false,
    tabsize=2,
    captionpos=b
}

\pagestyle{fancy}
\fancyhf{}
\fancyhead[L]{\small\textit{IPSL -- GeIT1, Base de Donn\'ees NoSQL}}
\fancyhead[R]{\small\textit{Mamadou Sow, El Hadji Ndao}}
\fancyfoot[C]{\thepage}
\renewcommand{\headrulewidth}{0.4pt}
\renewcommand{\footrulewidth}{0.4pt}

\graphicspath{{images/}}

\begin{document}

% ===================== PAGE DE GARDE =====================
\begin{titlepage}
\thispagestyle{empty}
\begin{center}

\vspace*{1.5cm}

Universit\'e Gaston Berger de Saint-Louis\\[0.3cm]
\textbf{\textit{Institut Polytechnique de Saint-Louis}}\\[0.3cm]
Ann\'ee Universitaire 2025/2026

\vspace{1cm}
\noindent\rule{\textwidth}{0.5pt}

\begin{flushleft}
\textbf{Fili\`ere :} GeIT1 -- IPSL\\
\textbf{Mati\`ere :} Base de Donn\'ees NoSQL\\
\textbf{Semestre :} 2\\
\textbf{Fait le :} 18 Juillet 2026
\end{flushleft}
\begin{flushright}
\vspace{-2.5cm}
\textbf{Professeur :} \textit{Dr. Pape Samour Diop}\\
\textbf{\'Etudiants :} \textit{Mamadou Sow}\\
\phantom{\textbf{\'Etudiants :}} \textit{El Hadji Ndao}
\end{flushright}

\noindent\rule{\textwidth}{0.5pt}

\vspace{3cm}

{\LARGE \textbf{Rapport de Projet}}

\vspace{1cm}

{\Huge \color{blue!70!black} Syst\`eme de Gestion de Files d'Attente}

\vspace{0.3cm}
{\large avec}
\vspace{0.3cm}

{\Huge \color{blue!70!black} Redis}

\vspace{0.5cm}

{\large Application de Messagerie (Python/Docker)}

\vspace{4cm}

{\large\color{blue!50!black}\textit{<< Premier arriv\'e, premier lu. >>}}

\end{center}
\end{titlepage}

% ===================== TABLE DES MATIERES =====================
\newpage
\tableofcontents
\newpage

% ===================== 1. INTRODUCTION =====================
\section{Introduction}

\subsection{Objectif du projet}

Ce projet consiste \`a d\'evelopper un syst\`eme de messagerie instantan\'ee o\`u Redis est utilis\'e pour g\'erer les files d'attente de messages entre utilisateurs. Chaque utilisateur poss\`ede une file de messages (liste Redis) qu'il peut consulter, vider ou enrichir.

L'objectif g\'en\'eral est d'utiliser Redis pour g\'erer des files d'attente de mani\`ere rapide et efficace dans une application temps r\'eel.

\subsection{Objectifs sp\'ecifiques}

\begin{itemize}
    \item Utiliser les structures de donn\'ees Redis : \texttt{LPUSH}, \texttt{RPUSH}, \texttt{LRANGE}, \texttt{LPOP}, \texttt{RPOP}
    \item Impl\'ementer un syst\`eme de messagerie simple entre utilisateurs
    \item G\'erer la persistance et l'expiration des messages
    \item Simuler des \'echanges et visualiser les files d'attente
\end{itemize}

\subsection{Technologies utilis\'ees}

\begin{table}[H]
\centering
\begin{tabular}{|l|l|l|}
\hline
\textbf{Technologie} & \textbf{Version} & \textbf{R\^ole} \\
\hline
Redis & 7 & Base de donn\'ees en m\'emoire \\
Python & 3.14 & Logique applicative \\
Docker Compose & 3.8 & Conteneurisation de Redis \\
redis-py & 5.0.1 & Biblioth\`eque Python pour Redis \\
\hline
\end{tabular}
\caption{Technologies utilis\'ees}
\end{table}

% ===================== 2. MODELISATION =====================
\section{Mod\'elisation}

\subsection{Choix de la structure de donn\'ees}

Nous avons choisi les \textbf{Listes Redis} comme structure principale car elles offrent :
\begin{itemize}
    \item Un comportement natif de file d'attente (FIFO)
    \item Des op\'erations O(1) pour l'insertion (\texttt{LPUSH}) et la lecture (\texttt{RPOP})
    \item La possibilit\'e de consulter sans supprimer (\texttt{LRANGE})
\end{itemize}

\subsection{Convention de nommage des cl\'es}

\begin{table}[H]
\centering
\begin{tabular}{|l|l|l|}
\hline
\textbf{Cl\'e Redis} & \textbf{Type} & \textbf{Usage} \\
\hline
\texttt{queue:<utilisateur>} & Liste & File de messages \\
\texttt{stats:total\_messages} & String & Compteur global \\
\texttt{stats:messages\_envoyes:<user>} & String & Messages envoy\'es par user \\
\texttt{stats:messages\_recus:<user>} & String & Messages re\c{c}us par user \\
\hline
\end{tabular}
\caption{Convention de nommage des cl\'es Redis}
\end{table}

\subsection{Format des messages}

Chaque message est stock\'e sous forme de cha\^ine avec le format :

\begin{lstlisting}
<id_message>|<expediteur>|<timestamp>|<contenu>
\end{lstlisting}

Exemple :
\begin{lstlisting}
msg_1784407000469|momar|1784407000|Reunion d'equipe reporte a jeudi.
\end{lstlisting}

\subsection{Sch\'ema du flux de messages (FIFO)}

Le principe FIFO est assur\'e par :
\begin{itemize}
    \item \texttt{LPUSH} : insertion en t\^ete de liste (\`a gauche)
    \item \texttt{RPOP} : extraction en queue de liste (\`a droite)/ ou \texttt{LPOP} (a gauche)
\end{itemize}

Ainsi, le premier message envoy\'e est le premier lu (First In, First Out).



% ===================== 3. INSTALLATION =====================
\section{Installation et configuration}

\subsection{Docker Compose}

Le fichier \texttt{docker-compose.yml} d\'eploie une instance Redis 7 :

\begin{lstlisting}[language=bash,caption={docker-compose.yml}]
version: "3.8"
services:
  redis:
    image: redis:7
    container_name: redis-messagerie
    ports:
      - "6379:6379"
    volumes:
      - redis-data:/data
    command: ["redis-server", "--appendonly", "yes",
              "--appendfsync", "everysec"]
    restart: unless-stopped

volumes:
  redis-data:
\end{lstlisting}

\subsection{Configuration Redis}

\begin{itemize}
    \item \textbf{Persistance AOF} activ\'ee (\texttt{--appendonly yes}) : chaque \'ecriture est journalis\'ee
    \item \textbf{Synchronisation} : \texttt{appendfsync everysec} (compromis performance/s\'ecurit\'e)
    \item \textbf{Volume Docker} : les donn\'ees survivent aux red\'emarrages du conteneur
\end{itemize}

\subsection{Environnement Python}

\begin{lstlisting}[language=bash]
python3 -m venv venv
source venv/bin/activate.fish
pip install -r requirements.txt   # redis==5.0.1
\end{lstlisting}

% ===================== 4. IMPLEMENTATION =====================
\section{Impl\'ementation}

\subsection{Architecture du projet}

\begin{lstlisting}
projet9-redis-messagerie/
|-- app/
|   |-- __init__.py         # Package Python
|   |-- config.py           # Connexion Redis
|   |-- messagerie.py       # Fonctions de gestion des files
|   |-- main.py             # Interface console interactive
|   +-- simulation.py       # Simulation automatique
|-- docker-compose.yml
|-- script_redis.sh
|-- requirements.txt
+-- README.md
\end{lstlisting}

\subsection{Module \texttt{messagerie.py} -- Fonctions principales}

\begin{table}[H]
\centering
\begin{tabular}{|l|l|l|}
\hline
\textbf{Fonction} & \textbf{Commande Redis} & \textbf{Description} \\
\hline
\texttt{envoyer\_message()} & LPUSH, INCR & Envoie un message \\
\texttt{lire\_messages()} & LRANGE / RPOP & Lit les messages \\
\texttt{lire\_message\_suivant()} & RPOP & Lit le plus ancien (FIFO) \\
\texttt{compter\_messages()} & LLEN & Compte les messages \\
\texttt{vider\_file()} & DEL & Supprime la file \\
\texttt{definir\_expiration()} & EXPIRE & D\'efinit un TTL \\
\texttt{obtenir\_ttl()} & TTL & Temps avant expiration \\
\texttt{lister\_utilisateurs()} & KEYS & Files actives \\
\texttt{obtenir\_statistiques()} & GET, LLEN, TTL & Vue d'ensemble \\
\hline
\end{tabular}
\caption{Fonctions du module messagerie.py}
\end{table}

\subsection{Interface interactive (\texttt{main.py})}

L'interface console offre 9 options :
\begin{enumerate}
    \item Envoyer un message
    \item Lire mes messages (sans supprimer)
    \item Lire le message suivant (supprime de la file)
    \item Vider ma file de messages
    \item D\'efinir un TTL sur ma file
    \item Voir le TTL de ma file
    \item Lister les utilisateurs actifs
    \item Statistiques globales
    \item Changer d'utilisateur
\end{enumerate}

\subsection{Simulateur (\texttt{simulation.py})}

Le simulateur teste le syst\`eme avec 3 utilisateurs (\texttt{mamadou}, \texttt{ldji}, \texttt{momar}) :
\begin{enumerate}
    \item Envoi al\'eatoire de 20 messages
    \item Lecture des messages de chaque utilisateur
    \item D\'emonstration FIFO
    \item D\'emonstration TTL (expiration)
    \item Mesure des performances (1000 messages)
    \item Statistiques globales
    \item R\'einitialisation de la base
\end{enumerate}

% ===================== 5. DEMONSTRATION =====================
\section{D\'emonstration}

\subsection{Interface interactive}

L'application se lance avec \texttt{python3 -m app.main}. L'utilisateur se connecte et acc\`ede au menu :

\begin{figure}[H]
\centering
\includegraphics[width=0.85\textwidth]{interface.png}
\caption{Connexion en tant que \texttt{mamadou} et menu principal}
\end{figure}

\subsection{Envoi de messages (\texttt{LPUSH})}

Les messages sont ins\'er\'es dans Redis avec la commande \texttt{LPUSH}. Le retour indique le nombre total d'\'el\'ements dans la liste apr\`es insertion :

\begin{figure}[H]
\centering
\includegraphics[width=0.9\textwidth]{envoi_sms.png}
\caption{LPUSH de messages dans les files via redis-cli}
\end{figure}

\subsection{Files actives (\texttt{KEYS})}

La commande \texttt{KEYS queue:*} affiche toutes les files existantes :

\begin{figure}[H]
\centering
\includegraphics[width=0.6\textwidth]{files_actives.png}
\caption{\texttt{KEYS queue:*} -- 3 files actives (mamadou, ldji, momar)}
\end{figure}

\subsection{Lecture des messages (\texttt{LRANGE})}

La commande \texttt{LRANGE queue:<user> 0 -1} affiche tous les messages d'une file sans les supprimer :

\begin{figure}[H]
\centering
\includegraphics[width=0.9\textwidth]{lire_sms.png}
\caption{Contenu des files de mamadou et ldji avec LRANGE}
\end{figure}

\subsection{Lecture FIFO (\texttt{RPOP})}

\texttt{RPOP} retire et retourne le message le plus ancien de la file (premier envoy\'e = premier lu) :

\begin{figure}[H]
\centering
\includegraphics[width=0.7\textwidth]{RPOP&LPOP.png}
\caption{RPOP retire le message le plus ancien -- principe FIFO}
\end{figure}

\subsection{Expiration TTL (\texttt{EXPIRE} / \texttt{TTL})}

On peut d\'efinir un d\'elai d'expiration sur une file. Apr\`es ce d\'elai, Redis supprime automatiquement la file et tous ses messages :

\begin{figure}[H]
\centering
\includegraphics[width=0.7\textwidth]{TTL.png}
\caption{EXPIRE queue:momar 30 -- TTL d\'ecroissant. TTL queue:mamadou = -1 (permanent)}
\end{figure}

\subsection{Suppression d'une file (\texttt{DEL})}

La commande \texttt{DEL} supprime imm\'ediatement une file et tous ses messages :

\begin{figure}[H]
\centering
\includegraphics[width=0.6\textwidth]{vider_file.png}
\caption{DEL queue:momar puis LLEN = 0 (file vide)}
\end{figure}

\subsection{Statistiques (\texttt{INCR} / \texttt{GET})}

Les compteurs sont incr\'ement\'es \`a chaque envoi avec \texttt{INCR} et consult\'es avec \texttt{GET} :

\begin{figure}[H]
\centering
\includegraphics[width=0.7\textwidth]{compteurs_Statistiques.png}
\caption{Compteurs de statistiques avec SET, INCR et GET}
\end{figure}

\subsection{Informations syst\`eme (\texttt{INFO})}

\begin{figure}[H]
\centering
\includegraphics[width=0.8\textwidth]{Info_Redis.png}
\caption{DBSIZE = 11 cl\'es, INFO memory (1.43 Mo utilis\'es)}
\end{figure}

\begin{figure}[H]
\centering
\includegraphics[width=0.6\textwidth]{keyspace.png}
\caption{INFO keyspace -- db0 : 11 cl\'es}
\end{figure}

\subsection{Simulation d'\'echanges automatiques}

Le simulateur g\'en\`ere 20 messages al\'eatoires entre mamadou, ldji et momar :

\begin{figure}[H]
\centering
\includegraphics[width=0.85\textwidth]{20_sms_aleatoires.png}
\caption{20 messages envoy\'es al\'eatoirement entre les 3 utilisateurs}
\end{figure}

\subsection{Monitoring temps r\'eel (\texttt{MONITOR})}

La commande \texttt{MONITOR} affiche toutes les commandes Redis en direct pendant que l'application tourne dans un autre terminal :

\begin{figure}[H]
\centering
\includegraphics[width=\textwidth]{Monitor.png}
\caption{MONITOR : commandes Redis en direct pendant la simulation Python}
\end{figure}

% ===================== 6. PERFORMANCES =====================
\section{Analyse des performances}

\subsection{R\'esultats mesur\'es}

Le test de performance (option 5 du simulateur) envoie et lit 1000 messages :

\begin{figure}[H]
\centering
\includegraphics[width=0.7\textwidth]{performances.png}
\caption{Mesure des performances -- envoi et lecture de 1000 messages}
\end{figure}

\begin{table}[H]
\centering
\begin{tabular}{|l|c|c|}
\hline
\textbf{Op\'eration} & \textbf{Dur\'ee} & \textbf{D\'ebit} \\
\hline
Envoi de 1000 messages (LPUSH) & 0.326s & \textbf{3 065 msg/s} \\
Lecture de 1000 messages (LRANGE) & 0.003s & \textbf{320 813 msg/s} \\
Lecture FIFO de 1000 messages (RPOP) & 0.086s & \textbf{11 589 msg/s} \\
\hline
\end{tabular}
\caption{R\'esultats du test de performances}
\end{table}

\subsection{Complexit\'e des op\'erations}

\begin{table}[H]
\centering
\begin{tabular}{|l|l|c|l|}
\hline
\textbf{Op\'eration} & \textbf{Commande} & \textbf{Complexit\'e} & \textbf{Explication} \\
\hline
Envoi & LPUSH & O(1) & Temps constant \\
Lecture FIFO & RPOP & O(1) & Temps constant \\
Lecture compl\`ete & LRANGE & O(N) & Lin\'eaire \\
Comptage & LLEN & O(1) & M\'etadonn\'ees \\
Expiration & EXPIRE & O(1) & Timer interne \\
Suppression & DEL & O(1) & Suppression cl\'e \\
\hline
\end{tabular}
\caption{Complexit\'e des op\'erations Redis utilis\'ees}
\end{table}

\subsection{Analyse des r\'esultats}

\begin{itemize}
    \item \textbf{LPUSH (3 065 msg/s)} : le d\'ebit est limit\'e par le r\'eseau Python $\rightarrow$ Redis (TCP localhost) et les 3 commandes \texttt{INCR} suppl\'ementaires par message pour les statistiques.
    \item \textbf{LRANGE (320 813 msg/s)} : extr\^emement rapide car c'est une seule commande qui retourne tous les messages en une fois. Les donn\'ees sont en m\'emoire.
    \item \textbf{RPOP (11 589 msg/s)} : plus rapide que LPUSH car il n'y a pas d'INCR. Chaque RPOP est un aller-retour r\'eseau individuel.
\end{itemize}

\subsection{Avantages de Redis pour les files de messages}

\begin{enumerate}
    \item \textbf{Latence sub-milliseconde} : donn\'ees en m\'emoire, pas d'acc\`es disque
    \item \textbf{Op\'erations atomiques} : LPUSH/RPOP sont thread-safe, pas de verrou n\'ecessaire
    \item \textbf{TTL natif} : expiration automatique g\'er\'ee par Redis
    \item \textbf{Persistance configurable} : AOF garantit la durabilit\'e sans sacrifier la performance
    \item \textbf{Simplicit\'e} : pas de broker complexe (vs RabbitMQ/Kafka pour un cas simple)
\end{enumerate}

\subsection{Limites}

\begin{enumerate}
    \item \textbf{M\'emoire limit\'ee} : toutes les donn\'ees sont en RAM (1.43 Mo pour notre test)
    \item \textbf{Pas de routage avanc\'e} : pas de topics, fanout, ou acknowledgement natif
    \item \textbf{LRANGE O(N)} : la lecture de toute la file est lin\'eaire en taille
    \item \textbf{Single-threaded} : Redis n'utilise qu'un seul c\oe ur CPU pour les commandes
\end{enumerate}

% ===================== 7. PERSISTANCE =====================
\section{Gestion de la persistance}

\subsection{Configuration AOF}

Notre instance Redis est configur\'ee avec la persistance AOF (Append Only File) :

\begin{lstlisting}
redis-server --appendonly yes --appendfsync everysec
\end{lstlisting}

\begin{itemize}
    \item \texttt{appendonly yes} : journalise chaque commande d'\'ecriture
    \item \texttt{appendfsync everysec} : synchronise le fichier AOF chaque seconde
\end{itemize}

\subsection{Fonctionnement}

\begin{enumerate}
    \item Chaque commande d'\'ecriture (LPUSH, INCR, DEL, EXPIRE) est ajout\'ee au fichier AOF
    \item En cas de red\'emarrage, Redis rejoue le fichier AOF pour restaurer l'\'etat exact
    \item Le fichier est stock\'e dans le volume Docker \texttt{/data/appendonlydir/}
\end{enumerate}

\subsection{Compromis performance / durabilit\'e}

\begin{table}[H]
\centering
\begin{tabular}{|l|l|l|}
\hline
\textbf{Mode} & \textbf{Perte maximale} & \textbf{Performance} \\
\hline
\texttt{appendfsync always} & 0 message & Plus lente \\
\textbf{\texttt{appendfsync everysec}} & \textbf{$\sim$1 seconde} & \textbf{Bon compromis (notre choix)} \\
\texttt{appendfsync no} & Variable (OS) & Plus rapide mais risqu\'ee \\
\hline
\end{tabular}
\caption{Modes de synchronisation AOF}
\end{table}

\subsection{Sauvegarde et restauration}

\begin{lstlisting}[language=bash]
# Sauvegarde manuelle
docker exec -it redis-messagerie redis-cli BGSAVE

# Verifier les fichiers de persistance
docker exec -it redis-messagerie ls -la /data/

# Restauration automatique au redemarrage
docker compose restart
\end{lstlisting}

Les donn\'ees sont conserv\'ees gr\^ace au volume Docker \texttt{redis-data}.

% ===================== 8. CONCLUSION =====================
\section{Conclusion}

Ce projet nous a permis de mettre en \oe uvre un syst\`eme de messagerie fonctionnel bas\'e sur les listes Redis, exploitant le principe FIFO (LPUSH/RPOP) pour g\'erer les files d'attente entre utilisateurs. Les fonctionnalit\'es d'envoi, de lecture, d'expiration automatique (TTL) et de statistiques ont \'et\'e impl\'ement\'ees avec succ\`es, et les tests de performance confirment la rapidit\'e de Redis avec un d\'ebit de plus de 3~000 messages par seconde en \'ecriture et 320~000 en lecture. La persistance AOF garantit la durabilit\'e des donn\'ees sans compromettre les performances. Le syst\`eme reste cependant limit\'e \`a une utilisation en console sans authentification ni notifications en temps r\'eel, et pourrait \'evoluer vers une interface web avec Redis Pub/Sub ou Streams pour une exp\'erience plus compl\`ete.

\end{document}
